
\section{Queueing Model and Problem Formulation}
\label{sec:model}

We formulate the queueing model, define the routing problem, and introduce the four GibbsQ policy classes.

\subsection{System Model}
\label{sec:model:system}

We adopt the standard parallel-server queueing model~\cite{kleinrock1975queueing,kelly2011reversibility,robert2003stochastic}.

\begin{definition}[Parallel heterogeneous queueing system]
\label{def:system}
Consider $\Nservers \geq 1$ parallel servers indexed by $i \in \{1, \dots, \Nservers\}$. Jobs arrive according to a Poisson process with rate $\arrate > 0$. Server~$i$ processes jobs independently at exponential rate $\srvrate_i > 0$. The queue-length process is denoted
\begin{equation}
\label{eq:queue_state}
\Qstate(t) = \bigl(\Qlen{1}(t), \dots, \Qlen{\Nservers}(t)\bigr) \in \Nats^\Nservers,
\end{equation}
where $\Nats = \{0, 1, 2, \dots\}$ is the set of non-negative integers. The state space is $\mathcal{X} = \Nats^\Nservers$.
\end{definition}

\begin{definition}[Total service capacity and load condition]
\label{def:load}
The total service capacity is
\begin{equation}
\label{eq:total_capacity}
\totcap \coloneqq \sum_{i=1}^{\Nservers} \srvrate_i.
\end{equation}
The \emph{load factor} is $\rho := \arrate/\totcap \in (0,1)$.
The system operates under the \emph{strict load condition}
\begin{equation}
\label{eq:load_condition}
\loadcond \quad (\text{equivalently, } \rho < 1).
\end{equation}
\end{definition}

\subsection{Routing Problem}
\label{sec:model:routing}

Upon each job arrival, a \emph{routing policy} selects the destination server. Formally, a routing policy is a measurable map
\begin{equation}
\label{eq:routing_map}
p \colon \Nats^\Nservers \to \Simplex,
\end{equation}
where $\Simplex = \bigl\{q \in [0,1]^\Nservers : \sum_{i=1}^\Nservers q_i = 1\bigr\}$ is the $(\Nservers-1)$-dimensional probability simplex; under policy~$p$, server~$i$ receives the arriving job with probability $\rprob{i}(\Qstate)$.

The primary performance objective is to minimize the steady-state expected total queue length:
\begin{equation}
\label{eq:objective}
\text{minimize} \quad \E\Bigl[\sum_{i=1}^{\Nservers} \Qlen{i}\Bigr].
\end{equation}

\subsection{Policy Classes}
\label{sec:model:policies}

The GibbsQ framework defines four policy classes: two analytical (provably stable) and two empirical.

\begin{definition}[Raw softmax routing]
\label{def:raw_softmax}
For temperature parameter $\temp > 0$, the raw softmax routing law assigns
\begin{equation}
\label{eq:raw_softmax}
\rprob{i}(\Qstate) = \frac{\exp(-\temp \Qlen{i})}{\sum_{j=1}^{\Nservers} \exp(-\temp \Qlen{j})}.
\end{equation}
This is a homogeneous entropy-regularized policy that ignores service-rate heterogeneity. Its stability is established in \cref{sec:softmax_theorem} via a Foster--Lyapunov argument.
\end{definition}


\begin{definition}[Unified Archimedean Softmax (UAS)]
\label{def:uas}
For temperature $\temp > 0$ and service rates $\srvrate_1, \dots, \srvrate_\Nservers > 0$, the UAS routing law assigns
\begin{equation}
\label{eq:uas_routing}
\rprob{i}(\Qstate, \srvrate) = \frac{\srvrate_i \exp\!\bigl(-\temp (\Qlen{i}+1)/\srvrate_i\bigr)}{\sum_{j=1}^{\Nservers} \srvrate_j \exp\!\bigl(-\temp (\Qlen{j}+1)/\srvrate_j\bigr)}.
\end{equation}
This is the heterogeneous extension of raw softmax that incorporates service rates into both the energy function and the prior. Its stability is established in \cref{sec:uas_theorem} via a weighted Lyapunov argument.
\end{definition}


\begin{definition}[Calibrated UAS]
\label{def:calibrated_uas}
The Calibrated UAS family generalizes UAS through three free parameters $(\calbeta, \calgamma, \caloffset)$:
\begin{equation}
\label{eq:calibrated_uas}
\rprob{i}^{\caluas}(\Qstate, \srvrate) \propto \srvrate_i^{\calgamma} \exp\!\left(-\temp \frac{\Qlen{i}+\caloffset}{\srvrate_i^{\calbeta}}\right),
\end{equation}
with benchmark parameters $(\calbeta, \calgamma, \caloffset) = (0.85, 0.5, 0.5)$. This family contains UAS as the special case $(\calbeta, \calgamma, \caloffset) = (1, 1, 1)$. Calibrated UAS is the full \emph{empirical} three-parameter family; formal stability claims are reserved for the named subclass introduced below and analyzed in \cref{sec:calibrated_uas}.
\end{definition}

\begin{definition}[Stability-Certified Calibrated UAS (SCUAS)]
\label{def:scuas}
For fixed parameters $(\calbeta, \calgamma, \caloffset)$, define the SCUAS routing law by
\begin{equation}
\label{eq:scuas_model}
\rprob{i}^{\scuas}(\Qstate, \srvrate) \propto \srvrate_i^{\calgamma} \exp\!\left(-\temp \frac{\Qlen{i}+\caloffset}{\srvrate_i^{\calbeta}}\right).
\end{equation}
SCUAS uses the same closed form as \eqref{eq:calibrated_uas}, but theorem-backed claims are made only for parameter settings satisfying the explicit sufficient drift condition in Theorem~\ref{thm:scuas}.
\end{definition}


\begin{definition}[N-GibbsQ]
\label{def:ngibbsq}
The $\ngibbsq$ policy is a learned neural routing policy parameterized by weights $\theta$:
\begin{equation}
\label{eq:ngibbsq}
\rprob{i}(\Qstate, \srvrate; \theta) = \bigl[\softmax\!\bigl(f_\theta(\Qstate, \srvrate)\bigr)\bigr]_i,
\end{equation}
where $f_\theta$ is a feedforward neural network. Training proceeds in two phases: behavior cloning ($\bc$) from an analytical teacher policy, followed by $\reinforce$ fine-tuning. The $\ngibbsq$ policy is an \emph{empirical} approximation layer and is \textbf{not} part of the theorem statements; see \cref{sec:neural_policy}.
\end{definition}

\subsection{Stability Framework}
\label{sec:model:stability}

The stability results use the continuous-time Foster--Lyapunov criterion from the Meyn--Tweedie framework~\cite{meyn2012markov,asmussen2003applied}.

A CTMC is \emph{positive Harris recurrent} if it possesses a unique stationary distribution $\pi$ and returns to every measurable set of positive $\pi$-measure in finite expected time~\cite{meyn2012markov}.
Positive Harris recurrence ensures that steady-state performance metrics, including $\E[\sum_i \Qlen{i}]$, are well-defined and can be estimated via ergodic averages.
The Foster--Lyapunov criterion provides a constructive sufficient condition:

\begin{definition}[Foster--Lyapunov drift condition]
\label{def:foster_lyapunov}
Let $\Qstate(t)$ be a continuous-time Markov chain (CTMC) on $\Nats^\Nservers$ with generator $\gen$. The chain satisfies the \emph{Foster--Lyapunov drift condition} if there exist a norm-like function $\lyap \colon \Nats^\Nservers \to [0,\infty)$, constants $\varepsilon > 0$ and $R < \infty$, and a finite set $C \subset \Nats^\Nservers$ such that
\begin{equation}
\label{eq:foster_lyapunov}
\gen \lyap(x) \leq -\varepsilon \abs{x}_1 + R \quad \text{for all } x \in \Nats^\Nservers.
\end{equation}
If this condition holds, then the CTMC is positive Harris recurrent.
\end{definition}

The raw softmax theorem (\cref{sec:softmax_theorem}), the UAS theorem (\cref{sec:uas_theorem}), and the SCUAS theorem (\cref{thm:scuas}) each verify \cref{eq:foster_lyapunov} with explicit Lyapunov functions and closed-form constants~$\varepsilon$ and~$R$.

\subsection{Notation Summary}
\label{sec:model:notation}

\cref{tab:notation} collects the principal notation used throughout the paper.

\begin{table}[htbp]
\centering
\small
\caption{Summary of principal notation. Symbols in the final block are introduced in Sections~\ref{sec:uas_theorem} and~\ref{sec:calibrated_uas} and appear here for reference.}
\label{tab:notation}
\renewcommand{\arraystretch}{1.30}
\begin{tabular}{@{\hspace{6pt}}l @{\hspace{14pt}} l@{\hspace{6pt}}}
\toprule
\textbf{Symbol} & \textbf{Description} \\
\midrule
$\Nservers$ & Number of parallel servers \\
$\arrate$ & Job arrival rate (Poisson) \\
$\srvrate_i$ & Service rate of server $i$ (exponential) \\
$\totcap = \sum_i \srvrate_i$ & Total service capacity \\
\addlinespace[3pt]
$\Qlen{i}(t)$ & Queue length at server $i$ at time $t$ \\
$\Qstate(t)$ & Queue-length vector $(\Qlen{1}(t), \dots, \Qlen{\Nservers}(t))$ \\
$\Simplex$ & Probability simplex in $\Reals^\Nservers$ \\
$\rprob{i}(\Qstate)$ & Routing probability for server $i$ \\
\addlinespace[3pt]
$\temp$ & Temperature (entropy-regularization strength) \\
$\entropy(p)$ & Shannon entropy $-\sum_i p_i \log p_i$ \\
$\KL{p}{r}$ & Kullback--Leibler divergence \\
\addlinespace[3pt]
$\gen$ & CTMC generator \\
$\lyap(\Qstate)$ & Lyapunov function \\
$\gen \lyap(\Qstate)$ & Generator drift \\
\addlinespace[3pt]
$\uasprior{i} = \srvrate_i / \totcap$ & UAS service-rate prior \\
$\uasenergy{i} = (\Qlen{i}+1)/\srvrate_i$ & UAS energy variable \\
$(\calbeta, \calgamma, \caloffset)$ & Calibrated UAS / SCUAS parameters \\
$\varepsilon_{\calbeta,\calgamma,\caloffset}$ & SCUAS drift-rate constant \\
\bottomrule
\end{tabular}
\end{table}

